\fsection{Ej 1 pág. 146}
\fsubsection
[\textcolor{waveBlue2}{Utiliza las TIC}]
{\textcolor{waveBlue2}{\boldit{Utiliza las TIC}}}
{
	Entra en la página web de \textsc{Cepsa} para ver su plan de sostenibilidad:
	\textless\,\url{https://bit.ly/3SIud71}\,\textgreater.
	Destaca alguna de sus líneas prioritarias de acción.
}

\begin{solution}

	\textbf{Plan de Sostenibilidad de Cepsa: \textit{``Driving Positive Impact''}}

	El plan de sostenibilidad de Cepsa se enmarca dentro de su estrategia corporativa para 2030,
	denominada \textit{Positive Motion}, con una inversión prevista de entre 7.000 y 8.000 millones
	de euros en la década. El objetivo central es que más de la mitad del EBITDA de la compañía
	provenga de negocios sostenibles en 2030.

	\medskip

	\textbf{Líneas prioritarias de acción:}

	\begin{itemize}
		\item \textbf{Descarbonización y cero emisiones netas:} Cepsa se compromete a alcanzar
		      la neutralidad de carbono en 2050, con reducciones progresivas de emisiones a lo largo
		      de la década. Este objetivo ha sido reconocido por agencias como Moody's, Sustainalytics
		      y S\&P Global.

		\item \textbf{Liderazgo en energías renovables y movilidad sostenible:} La compañía
		      apuesta por convertirse en referente del hidrógeno verde y los biocombustibles de segunda
		      generación en España y Portugal, así como en la electrificación del transporte mediante
		      una red de puntos de recarga.

		\item \textbf{Biocombustibles de segunda generación:} En colaboración con Bio-Oils, Cepsa
		      construye una planta en Palos de la Frontera (Huelva) capaz de procesar 600.000 toneladas
		      anuales de residuos orgánicos (aceites usados, desechos agrícolas) para producir hasta
		      500.000 toneladas de combustible sostenible de aviación (SAF) y diésel renovable (HVO).

		\item \textbf{Economía circular:} La empresa se fija como meta aumentar la circularidad
		      de sus residuos en un 50\,\% para 2030, y ya está explorando la producción de hidrógeno
		      verde mediante el reciclaje de aguas residuales.

		\item \textbf{Gestión hídrica responsable:} Cepsa se compromete a reducir la extracción
		      de agua dulce en zonas con estrés hídrico en un 20\,\% para 2025.

		\item \textbf{Diversidad e impacto social:} El plan incluye objetivos de diversidad de
		      género y programas de transformación cultural interna (programa \textit{Orion}) para
		      involucrar a toda la organización en la transición energética.
	\end{itemize}
	\medskip
\end{solution}
